
\PassOptionsToPackage{unicode}{hyperref}
\PassOptionsToPackage{hyphens}{url}
%
\documentclass[
]{article}
\usepackage{lmodern}
\usepackage{amssymb,amsmath}
\usepackage{ifxetex,ifluatex}
\usepackage{longtable}
\ifnum 0\ifxetex 1\fi\ifluatex 1\fi=0 % if pdftex
  \usepackage[T1]{fontenc}
  \usepackage[utf8]{inputenc}
  \usepackage{textcomp} % provide euro and other symbols
\else % if luatex or xetex
  \usepackage{unicode-math}
  \defaultfontfeatures{Scale=MatchLowercase}
  \defaultfontfeatures[\rmfamily]{Ligatures=TeX,Scale=1}
\fi
% Use upquote if available, for straight quotes in verbatim environments
\IfFileExists{upquote.sty}{\usepackage{upquote}}{}
\IfFileExists{microtype.sty}{% use microtype if available
  \usepackage[]{microtype}
  \UseMicrotypeSet[protrusion]{basicmath} % disable protrusion for tt fonts
}{}
\makeatletter
\@ifundefined{KOMAClassName}{% if non-KOMA class
  \IfFileExists{parskip.sty}{%
    \usepackage{parskip}
  }{% else
    \setlength{\parindent}{0pt}
    \setlength{\parskip}{6pt plus 2pt minus 1pt}}
}{% if KOMA class
  \KOMAoptions{parskip=half}}
\makeatother
\usepackage{xcolor}
\IfFileExists{xurl.sty}{\usepackage{xurl}}{} % add URL line breaks if available
\IfFileExists{bookmark.sty}{\usepackage{bookmark}}{\usepackage{hyperref}}
\hypersetup{
  hidelinks,
  pdfcreator={LaTeX via pandoc}}
\urlstyle{same} % disable monospaced font for URLs
\usepackage{longtable,booktabs}
% Correct order of tables after \paragraph or \subparagraph
\usepackage{etoolbox}
\makeatletter
\patchcmd\longtable{\par}{\if@noskipsec\mbox{}\fi\par}{}{}
\makeatother
% Allow footnotes in longtable head/foot
\IfFileExists{footnotehyper.sty}{\usepackage{footnotehyper}}{\usepackage{footnote}}
\makesavenoteenv{longtable}
\setlength{\emergencystretch}{3em} % prevent overfull lines
\providecommand{\tightlist}{%
  \setlength{\itemsep}{0pt}\setlength{\parskip}{0pt}}
\setcounter{secnumdepth}{-\maxdimen} % remove section numbering


%% Sets the margins
\usepackage{geometry}
\makeatletter
\def\@biblabel#1{\hspace*{-\labelsep}}
\makeatother
\geometry{left=1in,right=1
in,top=1in,bottom=1.2in}
\newdimen\dummy
\dummy=\oddsidemargin
\addtolength{\dummy}{72pt}
\marginparwidth=.5\dummy
\marginparsep=.1\dummy

\definecolor{statagreen}{RGB}{34, 139, 34} 

\begin{document}

\title{\textbf{Read Me File for: \\
Managers and Productivity in Retail}}
\author{Robert D. Metcalfe, Alexandre B. Sollaci, and Chad Syverson \\
(Journal of Political Economy, 2026)}	
\date{}
\maketitle 
\thispagestyle{empty}


\subsection{Overview}

The Replication folder contains 4 sub-folders that include all of the required data and codes to replicate the paper. The codes are separated in the "CodeA" and "CodeB" folders, which perform the analysis for companies A and B, respectively. The "Data" folder contains the two source datasets described below; all auxiliary datasets will be saved here as well. Finally, the "Results" folder contains another two sub-folders: "Figures", where all figures are saved, and "Tables", where data required to construct the tables in the paper are saved.

The code in this replication package constructs the analysis from the data of two large retailers using R, Stata, and Python. Within the "CodeA" and "CodeB" folders, three main files generate all of the empirical results in the paper. Simulated results are run through the files in "CodeA/NAM Simulations". The replicator should expect the code to run for about 1 hour for each company.

\subsection{Data Availability}

Data from the two retailers was provided under the condition that companies remain anonymous. As such, the data used for this paper cannot be made publicly available. The JPE Data Editor has conducted plausibility checks concerning data provenance under an NDA with the authors.

The data should be preserved in the restricted-access location of the Journal of Political Economy for as long as necessary for the successful replication of the paper. It should be deleted after the replication is confirmed.

The authors commit to preserving the data and code for a period of at least five years following the publication of the paper and to provide reasonable assistance to requests
for clarification and replication.

\subsubsection{Statement of Rights}
I certify that the authors of the manuscript have legitimate access to and permission to use the data used in this manuscript.


\subsubsection{Dataset List}

The main data sources for the analysis are the data provided by each company. They are saved on the "Data" folder as Stata data files: "A\_dataset.dta" and "B\_dataset.dta". A full data dictionary with the variables included in each of these datasets is provided at the end of this document in Tables \ref{tab:data_dictionaryA} and \ref{tab:data_dictionaryB} at the end of this file.

Several other datasets are created for the remainder of the analysis. These are listed and described in the table below. Please note that we use the letter "X" do denote datasets that are available for both companies (e.g., "X\_dataset.dta" indicates that the datasets "A\_dataset.dta" and "B\_dataset.dta" exist for companies A and B, respectively).
\begin{table}[h]
    \begin{center}
    \caption{Datasets Used in Paper}\label{tab:placeholder}
    \begin{tabular}{lp{4.5cm}p{6cm}} \toprule \\
        Dataset & Origin & Description \\ \hline
        X\_dataset.dta & Source data & Dataset obtained from retail companies (cleaned) \\
        X\_connected.dta & Constructed; see file named: X\_estimate\_fe.R & Identifies all connected sets in data \\
        X\_manager\_fe.dta & Constructed; see file named: X\_estimate\_fe.R & Contains the initial manager fixed effect estimates \\
        X\_store\_fe.dta & Constructed; see file named: X\_estimate\_fe.R & Contains the initial store fixed effect estimates \\
        X\_manager\_gid.dta & Constructed; see file named: X\_estimate\_gfe.R & Contains the manager group id \\
        X\_manager\_gfe.dta & Constructed; see file named: X\_estimate\_gfe.R & Contains the manager grouped fixed effect estimates \\
        X\_store\_gid.dta & Constructed; see file named: X\_estimate\_gfe.R & Contains the store group id \\
        X\_store\_gfe.dta & Constructed; see file named: X\_estimate\_gfe.R & Contains the store grouped fixed effect estimates \\
        X\_dataset\_fe\_eb.dta & Constructed; see file named: X\_fixed\_effects.do & Matches the manager and store fixed effects to the source data and computes the EB-adjusted fixed effect estimates \\ 
        X\_dataset\_gfe.dta & Constructed; see file named: X\_grouped\_fixed\_effects.do & Matches the manager and store grouped fixed effects to the source data \\
        X\_event\_study.dta & Constructed; see file named: X\_event\_study.do & Prepares the "X\_dataset\_fe\_eb.dta" to run the event studies \\
        simulate\_match.dta & Constructed; see file named: selection\_based\_NAM.do & Keeps track of simulated manager-store matches; not used for empirical results \\
        \bottomrule
    \end{tabular}
    \end{center}
    \noindent \footnotesize Note: None of the datasets can be made publicly available. See below for a description of the code scripts.
\end{table}

\subsection{Computational Requirements}

There are no specific hardware requirements to run the codes provided, and most modern laptops should be able to complete them in about 1 hour (for each company). For the paper, codes were run on a Lenovo ThinkPad T14s laptop, with 32GB of installed RAM and a 2.00 GHz processor (AMD Ryzen AI 7 PRO 350 w/ Radeon 860M). The operating system was Windows 11 Enterprise. The replicator should be able to run scripts on Stata, R, and Python. The specific libraries and dependencies required in each software are listed below.

\begin{itemize}
    \item R (version 4.4.0) will require the following packages to be installed:
        \begin{itemize}
            \item haven
            \item rstudioapi
            \item doBy
            \item lfe
            \item Matrix (this is not explicitly called as a library, but is required for `lfe' to work properly)
        \end{itemize}
    \item Stata (version 19) will require the following packages to be installed (all available through ssc):
        \begin{itemize}
            \item reghdfe
            \item ftools
            \item did\_imputation
            \item event\_plot
            \item binscatter
            \item grc1leg2
            \item white\_tableau scheme (can be installed through: ssc install schemepack)
        \end{itemize}
    \item Python (version 3.11.11) will require the following libraries:
        \begin{itemize}
            \item numpy
            \item pandas
            \item scipy
            \item matplotlib
            \item os
        \end{itemize}
\end{itemize}

In addition to those, replication will also require the \textbf{ebayes.ado} and \textbf{fese\_fast.ado} programs written by Adam Sacarny (see \url{https://sacarny.com/programs/}). For convenience, these can be found within the "Stata codes" sub-folder in the "CodeA" and "CodeB" folders.


\subsection{Instructions to Replicators}

Each "CodeX" folder contains (for X = A, B):
\begin{itemize}
    \item A "Stata codes" folder
    \item Three main scripts
        \begin{itemize}
            \item X\_estimate\_fe.R
            \item X\_estimate\_gfe.R
            \item run\_stata\_codes.do
        \end{itemize}
\end{itemize}

The "CodeA" folder also contains a "NAM Simulations" folder, which has the scripts and data required to run the simulations for the NAM results:
\begin{itemize}
    \item selection\_based\_NAM.do
    \item simulate\_selection.py
    \item calibrate\_NAM.py
    \item NAM\_data.xlsx
\end{itemize}

\subsubsection{Running the Replication Codes}

To successfully replicate the paper, the user should
\begin{enumerate}
    \item Run the two scripts in R ("X\_estimate\_fe.R" and "X\_estimate\_gfe.R", X = A,B) to estimate the manager and store fixed effects and grouped fixed effects. 
    \item Run the master file "run\_stata\_codes.do". This will call on several other files contained in the "Stata codes" folder and will provide all of the remaining empirical results discussed in the paper.
    \item Run the scripts within the "CodeA/NAM Simulations" (no particular order required). This will provide the remaining results pertaining to the simulation exercises.
\end{enumerate}

\subsubsection{Description of the Code Scripts}

We provide a detailed description of all scripts and the results they provide below. As before, we use X as a placeholder for A or B in the names of scripts that are used for both companies A and B.
\begin{itemize}
    \item \textbf{X\_estimate\_fe.R} provides the initial estimates of the manager, store, and time fixed effects from equation 1 in the paper. It also adjusts the variance of the fixed effects for limited mobility bias and computes the variance/covariance shares in the four largest connected sets, as described in {\color{blue}Table 1}.
    
    \underline{Constructing Table 1}: This script will save the variance/covariance shares in the files "j-X\_cs\_covar.txt", for j = 1,...,4 representing the four largest connected sets (for example, "2-A\_cs\_covar.txt" has the variance/covariance shares for the second largest connected set in Company A). These files can be found the "Results/Tables" folder. 

    \textbf{Notes}: 
    
    (a) The number of observations within each connected set shown in {\color{blue}Table 1} is obtained after merging the "X\_store\_fe.dta" and "X\_manager\_fe.dta" datasets back into the original "X\_dataset.dta", and is done in the "X\_fixed\_effects.do" file (it will be displayed in the Stata window when run).

    (b) The estimation of the fixed effects involves a random number generator; a seed has been set on line 17.

    \item \textbf{X\_estimate\_gfe.R} clusters all managers and stores into groups based on observables and estimates the fixed effect of each group using equation 1. As above, it also adjusts the variance and covariance of estimates for limited mobility bias and saves them in the file "X\_covar\_group.txt" (in this case there is a single connected set). This file is again found in the "Results/Tables" folder.

    \textbf{Note}: 
    
    The estimation of the fixed effects involves a random number generator; a seed has been set on line 22.
\end{itemize}

\underline{\textbf{A comment on navigating the code files}}: The order of the results in the code files do not always follow the order they are presented in the paper. To facilitate the navigation between code and results, the "run\_stata\_code.do" file indicates which results are produced in each do file it calls. These will appear as a {\color{statagreen}commented code} right before the "do [file name]" command. 

For example, the code excerpt below indicates that the file "A\_summary\_stats.do" produces Figures A.1 and C.1, and Tables B.1, C.1, and D.1. 
\begin{verbatim}
// Figures A.1 and C.1; Tables B.1, C.1, and D.1
do "$code_dir/A_summary_stats.do"
\end{verbatim}

In addition, the section of code that produces each result is also identified with commented code. For example, within the "A\_summary\_stats.do" file, Figure A.1 is produced between lines 116 and 158. Line 115 is therefore a comment: {\color{statagreen}*** Figure A.1: plot pre-move trends ***}. A similar pattern holds for all Figures and Tables produced in by files below.

\begin{itemize}
    \item \textbf{X\_fixed\_effects.do} [in Stata codes folder] takes the initial estimates of the fixed effects (above) and adjusts them using the Empirical Bayes Adjustment. This is done with the help of the programs \textbf{ebayes.ado} and \textbf{fese\_fast.ado}.

    This script also
        \begin{itemize}
            \item Calculates the number of managers, stores, and overall number of observations in each connected set. These resuts are used in {\color{blue}Table 1} and discussed in {\color{blue}Section 3.2} (see the paragraph that starts with "Our data contains about [...]").
            \item Plots the distribution of fixed effects ({\color{blue}Figure 3}, {\color{blue}Figure A.2}, and {\color{blue}Figure A.3}) and the number of stores and managers in each connected set ({\color{blue}Figure A.4});
            \item Runs the out-of-sample test ({\color{blue}Figure 4});
            \item Performs the series of robustness checks discussed in {\color{blue}Appendix C.2}. This includes {\color{blue}Figure C.2} and {\color{blue}Figure C.3}, as well as the correlations (paragraph on local shocks) and the $R^2$ numbers (paragraphs on persistence and learning) in {\color{blue}Appendix C.2}.
            \item Estimates the full covariance matrix for manager and store fixed effects and performs the exercises discussed in the "Correlated Measurement Error" section of {\color{blue}Appendix E.2} (see the last paragraph for the calculated ratios).
        \end{itemize}
        
    \item \textbf{X\_grouped\_fixed\_effects.do} [in Stata codes folder] analyzes the grouped fixed effects estimates, producing the Figures in {\color{blue}Appendix E.1} ({\color{blue}Figure E.1}, {\color{blue}Figure E.2}, and {\color{blue}Figure E.3}).

    \item \textbf{X\_event\_study.do} [in Stata codes folder] runs all of the analysis related to the event studies. This includes {\color{blue}Figure 5} and {\color{blue}Figure 8} for Company A, {\color{blue}Figure 6} and {\color{blue}Figure 9} company B, and all Figures in {\color{blue}Appendix G}. 

    Using the estimates from the event studies, this script also produces the results required to construct {\color{blue}Table 4}, which are displayed in Stata's interface when the code is run. The replicator can find this information next to the headers "\#\#\# Table 4 (Panel A) \#\#\#" and "\#\#\# Table 4 (Panel B) \#\#\#". These will indicate results before and after the change in manager by the indicator "post\_change" (=0 for before change; =1 for after change). The point estimates named "b\_tb" for a top-to-bottom change and "b\_bt" for a bottom-to-top change. 

    In addition, the script analyzes what changes for managers and stores when a store switches managers, producing {\color{blue}Figure 7} and {\color{blue}Figure A.5}. The tenure gap mentioned in {\color{blue}Section 4.1} are also calculated here, and will be displayed in the Stata interface under the heaer "Average tenure change (section 4.1)".
    
    Finally, after breaking managers and stores into terciles based on their qualities and once again looking at the initial matches and moves, the script produces {\color{blue}Figure D.1}, as well as {\color{blue}Table 2} and {\color{blue}Table 3}. Data for both tables are saved in the "Results/Tables" folder (under "Table\_2\_X.dta" and "Table\_3\_X.dta").

    \item \textbf{X\_summary\_stats.do} [in Stata codes folder] calculates the summary statistics mentioned in the beginning of {\color{blue}section 2} in the paper. These will be displayed in the Stata interface under the heading "number of managers per store".
    
    The script also does some of the preliminary analyzes of the difference between "movers" and "non-movers", plotting {\color{blue}Figure A.1} and producing all the necessary numbers for {\color{blue}Table B.1} (displayed in Stata's interface under the heading "\#\#\# Table B.1 \#\#\#"). 
    
    Next, it estimates how different variables impact the probability that a manager moves between stores, running the regressions reported in {\color{blue}Table D.1} (displayed in the Stata interface under the header "\#\#\# Table D.1 \#\#\#"). 
    
    Finally, this script compares different measures of productivity, running all of the required analysis for the discussion in {\color{blue}Appendix C.1} (including {\color{blue}Figure C.1} and {\color{blue}Table C.1}, displayed in Stata's interface under the header "\#\#\# Table C.1 \#\#\#")

    \item \textbf{X\_counterfactuals.do} [in Stata codes folder] does all of the analysis reported in {\color{blue}Appendix F}, including {\color{blue}Table F.1} (again displayed in Stata's interface under the header "\#\#\# Table F.1 \#\#\#").

    \item \textbf{B\_gender\_moves.do} [in Stata codes folder] (only available for company B) runs the regressions required for {\color{blue}Table B.2} (displayed in the Stata interface under the header "\#\#\# Table B.2 \#\#\#"). It also calculates the share of male and female movers (displayed in the Stata interface under the header "Prob(mover|female)") to arrive at the 7 percentage point difference mentioned in {\color{blue}Footnote 7}. 
\end{itemize}

Within the "CodeA/NAM Simulations" folder the replicator will find the following files:
\begin{itemize}
    \item \textbf{selection\_based\_NAM.do} is a do file that generates {\color{blue}Figure 2} and {\color{blue}Figure D.2}.

    \textbf{Note}:

    The simulation requires random draws from a normal distribution; a seed has been set in line 13.

    \item \textbf{simulate\_selection.py} is a python script that generalizes the simulations from the "selection\_based\_NAM.do" file into 20 rounds of (re)matching of managers and stores. This is used to produce the three charts in {\color{blue}Figure D.3}.

    \textbf{Note}:

    The simulation requires random draws from a normal distribution; a seed has been set in line 15.
    
    \item \textbf{calibrate\_NAM.py} uses the observed variances and covariances in the data ("NAM\_data.xlsx") to recover the size of the bias in the correlation between managers and stores using the expression derived in the "\textbf{Discussion and Extensions}" part of {\color{blue}Section 3.3} and further developed in {\color{blue}Appendix D.2.1}.
    
    \item \textbf{NAM\_data.xlsx} is an excel file that simply calculates the weighted averages of covariances from {\color{blue}Table B.1}. The numbers found in this file are used in the 'calibrate\_NAM.py' program.
\end{itemize}










\clearpage 
\appendix

\begin{longtable}{ll} \caption{Data Dictionary, Company A} \label{tab:data_dictionaryA} \\ \toprule
Variable Name & Description \\ \hline 
id & (manager,store) pair id \\
store\_id & store id (numeric) \\
manager\_id & manager id (numeric) \\
time & period id \\
year & year of observations \\
month & month of observation (name)  \\
month\_num & month of observation (numeric) \\
location & name of store's location \\
location\_num & store location (numeric)  \\
state & US State of store (name) \\
state\_num & US State of store (numeric)  \\
county & County of store (name)  \\
countyfip & County FIP  \\
statefip & State FIP  \\
postal\_code & postal code  \\
timezone & time zone of store \\
remodel & indicator if store was remodeled  \\
(comp name)startdate & date that mng started in company \\
positionstartdate & date that mng started position \\
job\_description & job title  \\
paystatus & pay status  \\
supervisor\_id &  \\
supv\_title & supervisor title  \\
pos\_start\_month & month that mng started position  \\
pos\_start\_year & year that mng started position  \\
(comp name)\_start\_month & month that mng started in company  \\
(comp name)\_start\_year & year that mng started in company  \\
manager\_tenure\_pos & mng tenure in position  \\
manager\_tenure\_comp & mng tenure in company  \\
sales & total sales in month  \\
labor\_cost & cost of labor  \\
manager\_salary & manager's salary  \\
pt\_count & number of part time employees \\
ft\_count & numer of full-time employees  \\
oh\_cost\_consigned & overhead cost (consigned)  \\
oh\_cost\_owned & overhead cost (owned)  \\
oh\_cost\_total & overhead cost (total)  \\
sales\_area & size of sales area in store  \\
store\_area & size of store  \\
fte\_count & number of full-time equivalent employees  \\
log\_sales & log of sales  \\
log\_manager\_salary & log of manager salary  \\
log\_manager\_tenure\_pos & log of tenure in position  \\
log\_manager\_tenure\_comp & log of tenure in company  \\
log\_ft\_count & log of full-time employees  \\
log\_pt\_count & log of part-time employees  \\
log\_fte\_count & log of full-time equivalent employees  \\
log\_store\_area & log of store area  \\
log\_sales\_area & log of sales area  \\
capital\_cost & Cost of capital  \\
oh\_cost & overhead cost (same as oh\_cost\_total)  \\
log\_oh\_cost & log of overhead cost  \\
closest\_competitor & distance to closest competitor (miles)  \\
closest\_[competitor 1] & distance to closest competitor 1 (company's name in brackets)  \\
closest\_[competitor 2] & distance to closest competitor 2  \\
closest\_[competitor 3] & distance to closest competitor 3  \\
closest\_[competitor 4] & distance to closest competitor 4  \\
num\_mng\_2019 & number of different managers in 2019  \\
num\_close\_competitors & number of close competitors (within 1 mile)  \\
sun\_open\_hour & hour that store opens on Sunday  \\
sun\_close\_hour & hour that store closes on Sunday  \\
sun\_num\_hours & number of hours open on Sunday  \\
mon\_open\_hour & hour that store opens on Monday  \\
mon\_close\_hour & similar as above  \\
mon\_num\_hours &   \\
tue\_open\_hour &   \\
tue\_close\_hour &   \\
tue\_num\_hours &   \\
wed\_open\_hour &   \\
wed\_close\_hour &   \\
wed\_num\_hours &   \\
thu\_open\_hour &   \\
thu\_close\_hour &   \\
thu\_num\_hours &   \\
fri\_open\_hour &   \\
fri\_close\_hour &   \\
fri\_num\_hours &   \\
sat\_open\_hour &   \\
sat\_close\_hour &   \\
sat\_num\_hours &   \\
week\_num\_hours & number of hours open during week days  \\
weekend\_num\_hours & number of hours open during weekends  \\
sun\_log\_hours & log of hours open on Sunday  \\
mon\_log\_hours & log oh hours open on Monday  \\
tue\_log\_hours & similar to the above  \\
wed\_log\_hours &   \\
thu\_log\_hours &   \\
fri\_log\_hours &   \\
sat\_log\_hours &   \\
week\_log\_hours &   \\
weekend\_log\_hours &   \\ \bottomrule
\end{longtable}

\clearpage

\begin{longtable}{ll} \caption{Data Dictionary: Company B} \label{tab:data_dictionaryB} \\ \toprule
Variable Name & Description \\ \hline 
store\_id & store id (numeric) \\
manager\_id & manager id (numeric) \\
time & time period \\
year & year of observation \\
month & month of observation \\
elec\_consumption & consumption of electricity in month \\
elec\_cost & cost of electricity in month \\
gas\_consumption & consumption of gas in month \\
gas\_cost & cost of gas in month \\
competitor\_count & number of close competitors \\
mean\_dist\_km & mean distance to competitors in kilometers \\
mean\_drivetime & mean drivetime to competitors \\
mean\_dist\_drive\_km & mean driving distance to competitors in kilometers \\
min\_distance\_km & minimum distance to competitor in kilometers \\
min\_drivetime & minimum drive time to competitor \\
min\_dist\_drive\_km & minimum driving distance to competitor in kilometers \\
(competitor 1) & distance to (competitor 1) [name of competitor in parenthesis) \\
(competitor 2) & distance to (competitor 2) \\
(competitor 3) & distance to (competitor 3) \\
(competitor 4) & distance to (competitor 4) \\
(competitor 5) & distance to (competitor 5) \\
(competitor 6) & distance to (competitor 6) \\
(competitor 7) & distance to (competitor 7) \\
area\_total & size of store \\
area\_sales & size of sales area in store \\
store\_name & store's name \\
format & store format (flagship, local, others) \\
format\_num & store format (numeric) \\
type & store type (convenience, large, small, others) \\
opening\_date & store opening date \\
region & store region (numeric) \\
region\_name & store region name\\
area & store area (location, numeric) \\
area\_name & store area name \\
town & store town \\
county & store county \\
postcode & store postal code \\
longitude & store's location longitude \\
latitude & store's location latitude \\
country & store country \\
location\_code & store's location code \\
saturday\_opening & Saturday opening hour \\
saturday\_closing & Saturday closing hour \\
sunday\_opening & similar to above \\
sunday\_closing &  \\
weekday\_opening &  \\
weekday\_closing &  \\
weekday\_hours & number of minutes open during weekdays \\
saturday\_hours & number of minutes open during Saturdays \\
sunday\_hours & number of minutes open during Sundays \\
weekday\_hours\_total & total number of minutes open during weekdays \\
weekend\_hours & number of minutes open during the full weekend \\
week\_hours & total number of minutes open over the entire week (including weekends) \\
revenue\_total & total revenue during month \\
manager\_startdate & date when manager started job \\
months\_on\_job & number of months that manager has spent on job \\
hc\_count & number of employees \\
fte\_count & numer of full-time equivalent employees \\
elec\_use\_sqft & electricity consumption per squared foot in store \\
gas\_use\_sqft & gas consumption per squared foot in store \\
elec\_price & electricity price (constructed as cost/consumption) \\
gas\_price & gas price (constructed as cost/consumption) \\
energy\_cost & energy (electricty + gas) cost \\
log\_revenue\_total & log of total revenue \\
log\_energy\_cost & log of energy cost \\
log\_fte\_count & log of full-time equivalent employees \\
log\_hc\_count & log of number of employees \\
location\_id & store's city/location (numeric) \\
location\_name & store's city/location name \\
location\_type & location type (admin county, metropolitan, borough, etc.) \\
number\_active\_managers & number of active manager in store \\
multiple\_managers & indicator of multiple managers in store \\
number\_stores & number of store that manager in observation works in \\
multiple\_stores & indicator of manager that works in multiple stores \\
open\_year & year when store was opened \\
open\_month & month when store was opened \\
store\_age & store's age \\
log\_store\_age & log of store age \\
log\_store\_area & log of store size \\
log\_week\_hours & log of minutes open during weekdays \\
log\_weekend\_hours & log of minutes open during weekends \\
female & indicator of female manager \\
title & manager title (Mr., Mrs., Ms., others) \\
firstname & manager's first name \\
lastname & manager's last name \\
doj & date when manager joined company \\
join\_year & year when manager joined company \\
join\_month & month when manager joined company \\
manager\_tenure & manager's tenure in job/position \\
manager\_tenure\_start & manager's tenure in company \\
log\_manager\_tenure & log of manager tenure in position \\
log\_manager\_tenure\_start & log of manager tenure in company when they became a manager \\ \bottomrule
\end{longtable}



\end{document}


\hypertarget{template-readme-and-guidance}{%
\section{Template README and
Guidance}\label{template-readme-and-guidance}}

\begin{quote}
INSTRUCTIONS: This README suggests structure and content that have been
approved by various journals, see
\href{https://social-science-data-editors.github.io/template_README/Endorsers.html}{Endorsers}.
It is available as
\href{https://github.com/social-science-data-editors/template_README/blob/master/template-README.html}{Markdown/txt},
\href{https://social-science-data-editors.github.io/template_README/templates/README.docx}{Word},
\href{https://social-science-data-editors.github.io/template_README/templates/README.tex}{LaTeX},
and
\href{https://social-science-data-editors.github.io/template_README/templates/README.pdf}{PDF}.
In practice, there are many variations and complications, and authors
should feel free to adapt to their needs. All instructions can (should)
be removed from the final README (in Markdown, remove lines starting
with \texttt{\textgreater{}\ INSTRUCTIONS}). Please ensure that a PDF is
submitted in addition to the chosen native format. Please ensure that
the README is called ``README'' plus the appropriate suffix, not some
non-standard name. This helps replicators immediately locate the
necessary document.
\end{quote}

\hypertarget{overview}{%
\subsection{Overview}\label{overview}}

\begin{quote}
INSTRUCTIONS: The typical README in social science journals serves the
purpose of guiding a reader through the available material and a route
to replicating the results in the research paper. Start by providing a
brief overview of the available material and a brief guide as to how to
proceed from beginning to end.
\end{quote}

Example: The code in this replication package constructs the analysis
file from the three data sources (Ruggles et al, 2018; Inglehart et al,
2019; BEA, 2016) using Stata and Julia. Two main files run all of the
code to generate the data for the 15 figures and 3 tables in the paper.
The replicator should expect the code to run for about 14 hours.

\hypertarget{data-availability-and-provenance-statements}{%
\subsection{Data Availability and Provenance
Statements}\label{data-availability-and-provenance-statements}}

\begin{quote}
INSTRUCTIONS: Every README should contain a description of the origin
(provenance), location and accessibility (data availability) of the data
used in the article. These descriptions are generally referred to as
``Data Availability Statements'' (DAS). However, in some cases, there is
no external data used.
\end{quote}

\begin{itemize}
\tightlist
\item[$\square$]
  This paper does not involve analysis of external data (i.e., no data
  are used or the only data are generated by the authors via simulation
  in their code).
\end{itemize}

\begin{quote}
If box above is checked and if no simulated/synthetic data files are
provided by the authors, please skip directly to the section on
\protect\hyperlink{computational-requirements}{Computational
Requirements}. Otherwise, continue.
\end{quote}

\begin{quote}
INSTRUCTIONS: - When the authors are \textbf{secondary data users} (they
did not generate the data), the provenance and DAS coincide, and should
describe the condition under which (a) the current authors (b) any
future users might access the data. - When the data were generated (by
the authors) in the course of conducting (lab or field)
\textbf{experiments}, or were collected as part of \textbf{surveys},
then the description of the provenance should describe the data
generating process, i.e., survey or experimental procedures: -
Experiments: complete sets of experimental instructions, questionnaires,
stimuli for all conditions, potentially screenshots, scripts for
experimenters or research assistants, as well as for subject eligibility
criteria (e.g.~selection criteria, exclusions), recruitment waves,
demographics of subject pool used. - For lab experiments specifically, a
description of any pilot sessions/studies, and computer programs,
configuration files, or scripts used to run the experiment. - For
surveys, the whole questionnaire (code or images/PDF) including survey
logic if not linear, interviewer instructions, enumeration lists, sample
selection criteria.

The information should describe ALL data used, regardless of whether
they are provided as part of the replication archive or not, and
regardless of size or scope. The DAS should provide enough information
that a replicator can obtain the data from the original source, even if
the file is provided.

For instance, if using GDP deflators, the source of the deflators
(e.g.~at the national statistical office) should also be listed here. If
any of this information has been provided in a pre-registration, then a
link to that registration may (partially) suffice.

DAS can be complex and varied. Examples are provided
\href{https://social-science-data-editors.github.io/guidance/Requested_information_dcas.html}{here},
and below.

Importantly, if providing the data as part of the replication package,
authors should be clear about whether they have the \textbf{rights} to
distribute the data. Data may be subject to distribution restrictions
due to sensitivity, IRB, proprietary clauses in the data use agreement,
etc.

NOTE: DAS do not replace Data Citations (see
\href{https://social-science-data-editors.github.io/guidance/Data_citation_guidance.html}{Guidance}).
Rather, they augment them. Depending on journal requirements and to some
extent stylistic considerations, data citations should appear in the
main article, in an appendix, or in the README. However, data citations
only provide information \textbf{where} to find the data, not
\textbf{how to access} those data. Thus, DAS augment data citations by
going into additional detail that allow a researcher to assess cost,
complexity, and availability over time of the data used by the original
author.
\end{quote}

\hypertarget{statement-about-rights}{%
\subsubsection{Statement about Rights}\label{statement-about-rights}}

\begin{itemize}
\tightlist
\item[$\square$]
  I certify that the author(s) of the manuscript have legitimate access
  to and permission to use the data used in this manuscript.
\item[$\square$]
  I certify that the author(s) of the manuscript have documented
  permission to redistribute/publish the data contained within this
  replication package. Appropriate permission are documented in the
  \href{https://social-science-data-editors.github.io/template_README/LICENSE.txt}{LICENSE.txt}
  file.
\end{itemize}

\hypertarget{optional-but-recommended-license-for-data}{%
\subsubsection{(Optional, but recommended) License for
Data}\label{optional-but-recommended-license-for-data}}

\begin{quote}
INSTRUCTIONS: Most data repositories provide for a default license, but
do not impose a specific license. Authors should actively select a
license. This should be provided in a LICENSE.txt file, separately from
the README, possibly combined with the license for any code. Some data
may be subject to inherited license requirements, i.e., the data
provider may allow for redistribution only if the data is licensed under
specific rules - authors should check with their data providers. For
instance, a data use license might require that users - the current
author, but also any subsequent users - cite the data provider.
Licensing can be complex. Some non-legal guidance may be found
\href{https://social-science-data-editors.github.io/guidance/Licensing_guidance.html}{here}.
For multiple licenses within a data package, the \texttt{LICENSE.txt}
file might contain the concatenation of all the licenses that apply (for
instance, a custom license for one file, plus a CC-BY license for
another file).

NOTE: In many cases, it is not up to the creator of the replication
package to simply define a license, a license may be \emph{sticky} and
be defined by the original data creator.
\end{quote}

\emph{Example:} The data are licensed under a Creative Commons/CC-BY-NC
license. See LICENSE.txt for details.

\hypertarget{summary-of-availability}{%
\subsubsection{Summary of Availability}\label{summary-of-availability}}

\begin{itemize}
\tightlist
\item[$\square$]
  All data \textbf{are} publicly available.
\item[$\square$]
  Some data \textbf{cannot be made} publicly available.
\item[$\square$]
  \textbf{No data can be made} publicly available.
\end{itemize}

\begin{quote}
INSTRUCTIONS: If \emph{No or some data can be made} \textbf{publicly}
available, the journal may require that data be preserved for a
specified time.
\end{quote}

\begin{itemize}
\tightlist
\item[$\square$]
  Confidential data used in this paper and not provided as part of the
  public replication package will be preserved for \_\_\_ years after
  publication, in accordance with journal policies.
\end{itemize}

\hypertarget{details-on-each-data-source}{%
\subsubsection{Details on each Data
Source}\label{details-on-each-data-source}}

\begin{quote}
INSTRUCTIONS: For each data source, list the file that contains data
from that source here; if providing combined/derived datafiles, list
them separately after the DAS. For each data source or file, as
appropriate,

\begin{itemize}
\tightlist
\item
  Describe the format (open formats preferred, but some
  software-specific formats OK if open-source readers available):
  \texttt{.dta}, \texttt{.xlsx}, \texttt{.csv}, \texttt{netCDF}, etc.
\item
  Provide a data dictionairy, either as part of the archive (list the
  file name), or at a URL (list the URL). Some formats are
  self-describing \emph{if} they have the requisite information (e.g.,
  \texttt{.dta} should have both variable and value labels).
\item
  List availability within the package
\item
  Use proper bibliographic references in addition to a verbose
  description (and provide a bibliography at the end of the README,
  expanding those references)
\end{itemize}

A summary in tabular form can be useful:
\end{quote}

\begin{longtable}[]{@{}lllll@{}}
\toprule
\begin{minipage}[b]{0.17\columnwidth}\raggedright
Data.Name\strut
\end{minipage} & \begin{minipage}[b]{0.17\columnwidth}\raggedright
Data.Files\strut
\end{minipage} & \begin{minipage}[b]{0.17\columnwidth}\raggedright
Location\strut
\end{minipage} & \begin{minipage}[b]{0.17\columnwidth}\raggedright
Provided\strut
\end{minipage} & \begin{minipage}[b]{0.17\columnwidth}\raggedright
Citation\strut
\end{minipage}\tabularnewline
\midrule
\endhead
\begin{minipage}[t]{0.17\columnwidth}\raggedright
``Current Population Survey 2018''\strut
\end{minipage} & \begin{minipage}[t]{0.17\columnwidth}\raggedright
cepr\_march\_2018.dta\strut
\end{minipage} & \begin{minipage}[t]{0.17\columnwidth}\raggedright
data/\strut
\end{minipage} & \begin{minipage}[t]{0.17\columnwidth}\raggedright
TRUE\strut
\end{minipage} & \begin{minipage}[t]{0.17\columnwidth}\raggedright
CEPR (2018)\strut
\end{minipage}\tabularnewline
\begin{minipage}[t]{0.17\columnwidth}\raggedright
``Provincial Administration Reports''\strut
\end{minipage} & \begin{minipage}[t]{0.17\columnwidth}\raggedright
coast\_simplepoint2.csv; rivers\_simplepoint2.csv; RAIL\_dummies.dta;
railways\_Dissolve\_Simplify\_point2.csv\strut
\end{minipage} & \begin{minipage}[t]{0.17\columnwidth}\raggedright
Data/maps/\strut
\end{minipage} & \begin{minipage}[t]{0.17\columnwidth}\raggedright
TRUE\strut
\end{minipage} & \begin{minipage}[t]{0.17\columnwidth}\raggedright
Administration (2017)\strut
\end{minipage}\tabularnewline
\begin{minipage}[t]{0.17\columnwidth}\raggedright
``2017 SAT scores''\strut
\end{minipage} & \begin{minipage}[t]{0.17\columnwidth}\raggedright
Not available\strut
\end{minipage} & \begin{minipage}[t]{0.17\columnwidth}\raggedright
data/to\_clean/\strut
\end{minipage} & \begin{minipage}[t]{0.17\columnwidth}\raggedright
FALSE\strut
\end{minipage} & \begin{minipage}[t]{0.17\columnwidth}\raggedright
College Board (2020)\strut
\end{minipage}\tabularnewline
\bottomrule
\end{longtable}

where the \texttt{Data.Name} column is then expanded in the subsequent
paragraphs, and \texttt{CEPR\ (2018)} is resolved in the References
section of the README.

\hypertarget{example-for-public-use-data-collected-by-the-authors}{%
\subsubsection{Example for public use data collected by the
authors}\label{example-for-public-use-data-collected-by-the-authors}}

\begin{quote}
The {[}DATA TYPE{]} data used to support the findings of this study have
been deposited in the {[}NAME{]} repository ({[}DOI or OTHER PERSISTENT
IDENTIFIER{]}).
{[}\href{https://www.hindawi.com/research.data/\#statement.templates}{1}{]}.
The data were collected by the authors, and are available under a
Creative Commons Non-commercial license.
\end{quote}

\hypertarget{example-for-public-use-data-sourced-from-elsewhere-and-provided}{%
\subsubsection{Example for public use data sourced from elsewhere and
provided}\label{example-for-public-use-data-sourced-from-elsewhere-and-provided}}

\begin{quote}
Data on National Income and Product Accounts (NIPA) were downloaded from
the U.S. Bureau of Economic Analysis (BEA, 2016). We use Table 30. Data
can be downloaded from https://apps.bea.gov/regional/downloadzip.cfm,
under ``Personal Income (State and Local)'', select CAINC30: Economic
Profile by County, then download. Data can also be directly downloaded
using https://apps.bea.gov/regional/zip/CAINC30.zip. A copy of the data
is provided as part of this archive. The data are in the public domain.
\end{quote}

Datafile: \texttt{CAINC30\_\_ALL\_AREAS\_1969\_2018.csv}

\hypertarget{example-for-public-use-data-with-required-registration-and-provided-extract}{%
\subsubsection{Example for public use data with required registration
and provided
extract}\label{example-for-public-use-data-with-required-registration-and-provided-extract}}

\begin{quote}
The paper uses IPUMS Terra data (Ruggles et al, 2018). IPUMS-Terra does
not allow for redistribution, except for the purpose of replication
archives. Permissions as per https://terra.ipums.org/citation have been
obtained, and are documented within the ``data/IPUMS-terra'' folder.
\textgreater{} Note: the reference to ``Ruggles et al, 2018'' would be
resolved in the Reference section of this README, \textbf{and} in the
main manuscript.
\end{quote}

Datafile: \texttt{data/raw/ipums\_terra\_2018.dta}

\hypertarget{example-for-free-use-data-with-required-registration-extract-not-provided}{%
\subsubsection{Example for free use data with required registration,
extract not
provided}\label{example-for-free-use-data-with-required-registration-extract-not-provided}}

\begin{quote}
The paper uses data from the World Values Survey Wave 6 (Inglehart et
al, 2019). Data is subject to a redistribution restriction, but can be
freely downloaded from
http://www.worldvaluessurvey.org/WVSDocumentationWV6.jsp. Choose
\texttt{WV6\_Data\_Stata\_v20180912}, fill out the registration form,
including a brief description of the project, and agree to the
conditions of use. Note: ``the data files themselves are not
redistributed'' and other conditions. Save the file in the directory
\texttt{data/raw}.
\end{quote}

\begin{quote}
\begin{quote}
Note: the reference to ``Inglehart et al, 2018'' would be resolved in
the Reference section of this README, \textbf{and} in the main
manuscript.
\end{quote}
\end{quote}

Datafile: \texttt{data/raw/WV6\_Data\_Stata\_v20180912.dta} (not
provided)

\hypertarget{example-for-confidential-data}{%
\subsubsection{Example for confidential
data}\label{example-for-confidential-data}}

\begin{quote}
INSTRUCTIONS: Citing and describing confidential data, in particular
when it does not have a regular distribution channel or online landing
page, can be tricky. A citation can be crafted
(\href{https://social-science-data-editors.github.io/guidance/FAQ.html\#data-citation-without-online-link}{see
guidance}), and the DAS should describe how to access, whom to contact
(including the role of the particular person, should that person
retire), and other relevant information, such as required citizenship
status or cost.
\end{quote}

\begin{quote}
The data for this project (DESE, 2019) are confidential, but may be
obtained with Data Use Agreements with the Massachusetts Department of
Elementary and Secondary Education (DESE). Researchers interested in
access to the data may contact {[}NAME{]} at {[}EMAIL{]}, also see
www.doe.mass.edu/research/contact.html. It can take some months to
negotiate data use agreements and gain access to the data. The author
will assist with any reasonable replication attempts for two years
following publication.
\end{quote}

\hypertarget{example-for-confidential-census-bureau-data}{%
\subsubsection{Example for confidential Census Bureau
data}\label{example-for-confidential-census-bureau-data}}

\begin{quote}
All the results in the paper use confidential microdata from the U.S.
Census Bureau. To gain access to the Census microdata, follow the
directions here on how to write a proposal for access to the data via a
Federal Statistical Research Data Center:
https://www.census.gov/ces/rdcresearch/howtoapply.html. You must request
the following datasets in your proposal: 1. Longitudinal Business
Database (LBD), 2002 and 2007 2. Foreign Trade Database -- Import (IMP),
2002 and 2007 {[}\ldots{]}
\end{quote}

(adapted from \href{https://doi.org/10.1093/restud/rdw057}{Fort (2016)})

\hypertarget{example-for-preliminary-code-during-the-editorial-process}{%
\subsubsection{Example for preliminary code during the editorial
process}\label{example-for-preliminary-code-during-the-editorial-process}}

\begin{quote}
Code for data cleaning and analysis is provided as part of the
replication package. It is available at
https://dropbox.com/link/to/code/XYZ123ABC for review. It will be
uploaded to the {[}JOURNAL REPOSITORY{]} once the paper has been
conditionally accepted.
\end{quote}

\hypertarget{dataset-list}{%
\subsection{Dataset list}\label{dataset-list}}

\begin{quote}
INSTRUCTIONS: In some cases, authors will provide one dataset (file) per
data source, and the code to combine them. In others, in particular when
data access might be restrictive, the replication package may only
include derived/analysis data. Every file should be described. This can
be provided as a Excel/CSV table, or in the table below.
\end{quote}

\begin{quote}
INSTRUCTIONS: While it is often most convenient to provide data in the
native format of the software used to analyze and process the data, not
all formats are ``open'' and can be read by other (free) software. Data
should at a minimum be provided in formats that can be read by
open-source software (R, Python, others), and ideally be provided in
non-proprietary, archival-friendly formats.
\end{quote}

\begin{quote}
INSTRUCTIONS: All data files should be fully documented:
variables/columns should have labels (long-form meaningful names), and
values should be explained. This might mean generating a codebook,
pointing at a public codebook, or providing data in (non-proprietary)
formats that allow for a rich description. This is in particular
important for data that is not distributable.
\end{quote}

\begin{quote}
INSTRUCTIONS: Some journals require, and it is considered good practice,
to provide synthetic or simulated data that has some of the key
characteristics of the restricted-access data which are not provided.
The level of fidelity may vary - it may be useful for debugging only, or
it should allow to assess the key characteristics of the
statistical/econometric procedure or the main conclusions of the paper.
\end{quote}

\begin{longtable}[]{@{}llll@{}}
\toprule
\begin{minipage}[b]{0.26\columnwidth}\raggedright
Data file\strut
\end{minipage} & \begin{minipage}[b]{0.19\columnwidth}\raggedright
Source\strut
\end{minipage} & \begin{minipage}[b]{0.23\columnwidth}\raggedright
Notes\strut
\end{minipage} & \begin{minipage}[b]{0.21\columnwidth}\raggedright
Provided\strut
\end{minipage}\tabularnewline
\midrule
\endhead
\begin{minipage}[t]{0.26\columnwidth}\raggedright
\texttt{data/raw/lbd.dta}\strut
\end{minipage} & \begin{minipage}[t]{0.19\columnwidth}\raggedright
LBD\strut
\end{minipage} & \begin{minipage}[t]{0.23\columnwidth}\raggedright
Confidential\strut
\end{minipage} & \begin{minipage}[t]{0.21\columnwidth}\raggedright
No\strut
\end{minipage}\tabularnewline
\begin{minipage}[t]{0.26\columnwidth}\raggedright
\texttt{data/raw/terra.dta}\strut
\end{minipage} & \begin{minipage}[t]{0.19\columnwidth}\raggedright
IPUMS Terra\strut
\end{minipage} & \begin{minipage}[t]{0.23\columnwidth}\raggedright
As per terms of use\strut
\end{minipage} & \begin{minipage}[t]{0.21\columnwidth}\raggedright
Yes\strut
\end{minipage}\tabularnewline
\begin{minipage}[t]{0.26\columnwidth}\raggedright
\texttt{data/derived/regression\_input.dta}\strut
\end{minipage} & \begin{minipage}[t]{0.19\columnwidth}\raggedright
All listed\strut
\end{minipage} & \begin{minipage}[t]{0.23\columnwidth}\raggedright
Combines multiple data sources, serves as input for Table 2, 3 and
Figure 5.\strut
\end{minipage} & \begin{minipage}[t]{0.21\columnwidth}\raggedright
Yes\strut
\end{minipage}\tabularnewline
\bottomrule
\end{longtable}

\hypertarget{computational-requirements}{%
\subsection{Computational
requirements}\label{computational-requirements}}

\begin{quote}
INSTRUCTIONS: In general, the specific computer code used to generate
the results in the article will be within the repository that also
contains this README. However, other computational requirements - shared
libraries or code packages, required software, specific computing
hardware - may be important, and is always useful, for the goal of
replication. Some example text follows.
\end{quote}

\begin{quote}
INSTRUCTIONS: We strongly suggest providing setup scripts that
install/set up the environment. Sample scripts for
\href{https://github.com/gslab-econ/template/blob/master/config/config_stata.do}{Stata},
\href{https://github.com/labordynamicsinstitute/paper-template/blob/master/programs/global-libraries.R}{R},
\href{https://github.com/labordynamicsinstitute/paper-template/blob/master/programs/packages.jl}{Julia}
are easy to set up and implement. Specific software may have more
sophisticated tools:
\href{https://pip.pypa.io/en/stable/user_guide/\#ensuring-repeatability}{Python},
\href{https://julia.quantecon.org/more_julia/tools_editors.html\#Package-Environments}{Julia}.
\end{quote}

\hypertarget{software-requirements}{%
\subsubsection{Software Requirements}\label{software-requirements}}

\begin{quote}
INSTRUCTIONS: List all of the software requirements, up to and including
any operating system requirements, for the entire set of code. It is
suggested to distribute most dependencies together with the replication
package if allowed, in particular if sourced from unversioned code
repositories, Github repos, and personal webpages. In all cases, list
the version \emph{you} used. All packages should be listed in
human-readable form in this README, but should also be included in a
setup or install script.
\end{quote}

\begin{itemize}
\item[$\square$]
  The replication package contains one or more programs to install all
  dependencies and set up the necessary directory structure. {[}HIGHLY
  RECOMMENDED{]}
\item
  Stata (code was last run with version 15)

  \begin{itemize}
  \tightlist
  \item
    \texttt{estout} (as of 2018-05-12)
  \item
    \texttt{rdrobust} (as of 2019-01-05)
  \item
    the program ``\texttt{0\_setup.do}'' will install all dependencies
    locally, and should be run once.
  \end{itemize}
\item
  Python 3.6.4

  \begin{itemize}
  \tightlist
  \item
    \texttt{pandas} 0.24.2
  \item
    \texttt{numpy} 1.16.4
  \item
    the file ``\texttt{requirements.txt}'' lists these dependencies,
    please run ``\texttt{pip\ install\ -r\ requirements.txt}'' as the
    first step. See
    \url{https://pip.pypa.io/en/stable/user_guide/\#ensuring-repeatability}
    for further instructions on creating and using the
    ``\texttt{requirements.txt}'' file.
  \end{itemize}
\item
  Intel Fortran Compiler version 20200104
\item
  Matlab (code was run with Matlab Release 2018a)
\item
  R 3.4.3

  \begin{itemize}
  \tightlist
  \item
    \texttt{tidyr} (0.8.3)
  \item
    \texttt{rdrobust} (0.99.4)
  \item
    the file ``\texttt{0\_setup.R}'' will install all dependencies
    (latest version), and should be run once prior to running other
    programs.
  \end{itemize}
\end{itemize}

Portions of the code use bash scripting, which may require Linux.

Portions of the code use Powershell scripting, which may require Windows
10 or higher.

\hypertarget{controlled-randomness}{%
\subsubsection{Controlled Randomness}\label{controlled-randomness}}

\begin{quote}
INSTRUCTIONS: Some estimation code uses random numbers, almost always
provided by pseudorandom number generators (PRNGs). For reproducibility
purposes, these should be provided with a deterministic seed, so that
the sequence of numbers provided is the same for the original author and
any replicators. While this is not always possible, it is a requirement
by many journals' policies. The seed should be set once, and not use a
time-stamp. If using parallel processing, special care needs to be
taken. If using multiple programs in sequence, care must be taken on how
to call these programs, ideally from a main program, so that the
sequence is not altered. If no PRNG is used, check the other box.
\end{quote}

\begin{itemize}
\tightlist
\item[$\square$]
  Random seed is set at line \_\_\_\_\_ of program \_\_\_\_\_\_
\item[$\square$]
  No Pseudo random generator is used in the analysis described here.
\end{itemize}

\hypertarget{memory-runtime-storage-requirements}{%
\subsubsection{Memory, Runtime, Storage
Requirements}\label{memory-runtime-storage-requirements}}

\begin{quote}
INSTRUCTIONS: Memory and compute-time requirements may also be relevant
or even critical. Some example text follows. It may be useful to break
this out by Table/Figure/section of processing. For instance, some
estimation routines might run for weeks, but data prep and creating
figures might only take a few minutes. You should also describe how much
storage is required in addition to the space visible in the typical
repository, for instance, because data will be unzipped, data
downloaded, or temporary files written.
\end{quote}

\hypertarget{summary}{%
\paragraph{Summary}\label{summary}}

Approximate time needed to reproduce the analyses on a standard (CURRENT
YEAR) desktop machine:

\begin{itemize}
\tightlist
\item[$\square$]
  \textless10 minutes
\item[$\square$]
  10-60 minutes
\item[$\square$]
  1-2 hours
\item[$\square$]
  2-8 hours
\item[$\square$]
  8-24 hours
\item[$\square$]
  1-3 days
\item[$\square$]
  3-14 days
\item[$\square$]
  \textgreater{} 14 days
\end{itemize}

Approximate storage space needed:

\begin{itemize}
\item[$\square$]
  \textless{} 25 MBytes
\item[$\square$]
  25 MB - 250 MB
\item[$\square$]
  250 MB - 2 GB
\item[$\square$]
  2 GB - 25 GB
\item[$\square$]
  25 GB - 250 GB
\item[$\square$]
  \textgreater{} 250 GB
\item[$\square$]
  Not feasible to run on a desktop machine, as described below.
\end{itemize}

\hypertarget{details}{%
\paragraph{Details}\label{details}}

The code was last run on a \textbf{4-core Intel-based laptop with MacOS
version 10.14.4 with 200GB of free space}.

Portions of the code were last run on a \textbf{32-core Intel server
with 1024 GB of RAM, 12 TB of fast local storage}. Computation took
\textbf{734 hours}.

Portions of the code were last run on a \textbf{12-node AWS R3 cluster,
consuming 20,000 core-hours, with 2TB of attached storage}.

\begin{quote}
INSTRUCTIONS: Identifiying hardware and OS can be obtained through a
variety of ways: Some of these details can be found as follows:

\begin{itemize}
\tightlist
\item
  (Windows) by right-clicking on ``This PC'' in File Explorer and
  choosing ``Properties''
\item
  (Mac) Apple-menu \textgreater{} ``About this Mac''
\item
  (Linux) see code in
  \href{https://github.com/AEADataEditor/replication-template/blob/master/tools/linux-system-info.sh}{linux-system-info.sh}`
\end{itemize}
\end{quote}

\hypertarget{description-of-programscode}{%
\subsection{Description of
programs/code}\label{description-of-programscode}}

\begin{quote}
INSTRUCTIONS: Give a high-level overview of the program files and their
purpose. Remove redundant/ obsolete files from the Replication archive.
\end{quote}

\begin{itemize}
\tightlist
\item
  Programs in \texttt{programs/01\_dataprep} will extract and reformat
  all datasets referenced above. The file
  \texttt{programs/01\_dataprep/main.do} will run them all.
\item
  Programs in \texttt{programs/02\_analysis} generate all tables and
  figures in the main body of the article. The program
  \texttt{programs/02\_analysis/main.do} will run them all. Each program
  called from \texttt{main.do} identifies the table or figure it creates
  (e.g., \texttt{05\_table5.do}). Output files are called appropriate
  names (\texttt{table5.tex}, \texttt{figure12.png}) and should be easy
  to correlate with the manuscript.
\item
  Programs in \texttt{programs/03\_appendix} will generate all tables
  and figures in the online appendix. The program
  \texttt{programs/03\_appendix/main-appendix.do} will run them all.
\item
  Ado files have been stored in \texttt{programs/ado} and the
  \texttt{main.do} files set the ADO directories appropriately.
\item
  The program \texttt{programs/00\_setup.do} will populate the
  \texttt{programs/ado} directory with updated ado packages, but for
  purposes of exact reproduction, this is not needed. The file
  \texttt{programs/00\_setup.log} identifies the versions as they were
  last updated.
\item
  The program \texttt{programs/config.do} contains parameters used by
  all programs, including a random seed. Note that the random seed is
  set once for each of the two sequences (in \texttt{02\_analysis} and
  \texttt{03\_appendix}). If running in any order other than the one
  outlined below, your results may differ.
\end{itemize}

\hypertarget{optional-but-recommended-license-for-code}{%
\subsubsection{(Optional, but recommended) License for
Code}\label{optional-but-recommended-license-for-code}}

\begin{quote}
INSTRUCTIONS: Most journal repositories provide for a default license,
but do not impose a specific license. Authors should actively select a
license. This should be provided in a LICENSE.txt file, separately from
the README, possibly combined with the license for any data provided.
Some code may be subject to inherited license requirements, i.e., the
original code author may allow for redistribution only if the code is
licensed under specific rules - authors should check with their sources.
For instance, some code authors require that their article describing
the econometrics of the package be cited. Licensing can be complex. Some
non-legal guidance may be found
\href{https://social-science-data-editors.github.io/guidance/Licensing_guidance.html}{here}.
\end{quote}

The code is licensed under a MIT/BSD/GPL {[}choose one!{]} license. See
\href{https://social-science-data-editors.github.io/template_README/LICENSE.txt}{LICENSE.txt}
for details.

\hypertarget{instructions-to-replicators}{%
\subsection{Instructions to
Replicators}\label{instructions-to-replicators}}

\begin{quote}
INSTRUCTIONS: The first two sections ensure that the data and software
necessary to conduct the replication have been collected. This section
then describes a human-readable instruction to conduct the replication.
This may be simple, or may involve many complicated steps. It should be
a simple list, no excess prose. Strict linear sequence. If more than 4-5
manual steps, please wrap a main program/Makefile around them, in
logical sequences. Examples follow.
\end{quote}

\begin{itemize}
\tightlist
\item
  Edit \texttt{programs/config.do} to adjust the default path
\item
  Run \texttt{programs/00\_setup.do} once on a new system to set up the
  working environment.
\item
  Download the data files referenced above. Each should be stored in the
  prepared subdirectories of \texttt{data/}, in the format that you
  download them in. Do not unzip. Scripts are provided in each directory
  to download the public-use files. Confidential data files requested as
  part of your FSRDC project will appear in the \texttt{/data} folder.
  No further action is needed on the replicator's part.
\item
  Run \texttt{programs/01\_main.do} to run all steps in sequence.
\end{itemize}

\hypertarget{details-1}{%
\subsubsection{Details}\label{details-1}}

\begin{itemize}
\tightlist
\item
  \texttt{programs/00\_setup.do}: will create all output directories,
  install needed ado packages.

  \begin{itemize}
  \tightlist
  \item
    If wishing to update the ado packages used by this archive, change
    the parameter \texttt{update\_ado} to \texttt{yes}. However, this is
    not needed to successfully reproduce the manuscript tables.
  \end{itemize}
\item
  \texttt{programs/01\_dataprep}:

  \begin{itemize}
  \tightlist
  \item
    These programs were last run at various times in 2018.
  \item
    Order does not matter, all programs can be run in parallel, if
    needed.
  \item
    A \texttt{programs/01\_dataprep/main.do} will run them all in
    sequence, which should take about 2 hours.
  \end{itemize}
\item
  \texttt{programs/02\_analysis/main.do}.

  \begin{itemize}
  \tightlist
  \item
    If running programs individually, note that ORDER IS IMPORTANT.
  \item
    The programs were last run top to bottom on July 4, 2019.
  \end{itemize}
\item
  \texttt{programs/03\_appendix/main-appendix.do}. The programs were
  last run top to bottom on July 4, 2019.
\item
  Figure 1: The figure can be reproduced using the data provided in the
  folder ``2\_data/data\_map'', and ArcGIS Desktop (Version 10.7.1) by
  following these (manual) instructions:

  \begin{itemize}
  \tightlist
  \item
    Create a new map document in ArcGIS ArcMap, browse to the folder
    ``2\_data/data\_map'' in the ``Catalog'', with files
    ``provinceborders.shp'', ``lakes.shp'', and ``cities.shp''.
  \item
    Drop the files listed above onto the new map, creating three
    separate layers. Order them with ``lakes'' in the top layer and
    ``cities'' in the bottom layer.
  \item
    Right-click on the cities file, in properties choose the variable
    ``health''\ldots{} (more details)
  \end{itemize}
\end{itemize}

\hypertarget{list-of-tables-and-programs}{%
\subsection{List of tables and
programs}\label{list-of-tables-and-programs}}

\begin{quote}
INSTRUCTIONS: Your programs should clearly identify the tables and
figures as they appear in the manuscript, by number. Sometimes, this may
be obvious, e.g.~a program called ``\texttt{table1.do}'' generates a
file called \texttt{table1.png}. Sometimes, mnemonics are used, and a
mapping is necessary. In all circumstances, provide a list of tables and
figures, identifying the program (and possibly the line number) where a
figure is created.

NOTE: If the public repository is incomplete, because not all data can
be provided, as described in the data section, then the list of tables
should clearly indicate which tables, figures, and in-text numbers can
be reproduced with the public material provided.
\end{quote}

The provided code reproduces:

\begin{itemize}
\tightlist
\item[$\square$]
  All numbers provided in text in the paper
\item[$\square$]
  All tables and figures in the paper
\item[$\square$]
  Selected tables and figures in the paper, as explained and justified
  below.
\end{itemize}

\begin{longtable}[]{@{}lllll@{}}
\toprule
\begin{minipage}[b]{0.13\columnwidth}\raggedright
Figure/Table \#\strut
\end{minipage} & \begin{minipage}[b]{0.18\columnwidth}\raggedright
Program\strut
\end{minipage} & \begin{minipage}[b]{0.09\columnwidth}\raggedright
Line Number\strut
\end{minipage} & \begin{minipage}[b]{0.23\columnwidth}\raggedright
Output file\strut
\end{minipage} & \begin{minipage}[b]{0.23\columnwidth}\raggedright
Note\strut
\end{minipage}\tabularnewline
\midrule
\endhead
\begin{minipage}[t]{0.13\columnwidth}\raggedright
Table 1\strut
\end{minipage} & \begin{minipage}[t]{0.18\columnwidth}\raggedright
02\_analysis/table1.do\strut
\end{minipage} & \begin{minipage}[t]{0.09\columnwidth}\raggedright
\strut
\end{minipage} & \begin{minipage}[t]{0.23\columnwidth}\raggedright
summarystats.csv\strut
\end{minipage} & \begin{minipage}[t]{0.23\columnwidth}\raggedright
\strut
\end{minipage}\tabularnewline
\begin{minipage}[t]{0.13\columnwidth}\raggedright
Table 2\strut
\end{minipage} & \begin{minipage}[t]{0.18\columnwidth}\raggedright
02\_analysis/table2and3.do\strut
\end{minipage} & \begin{minipage}[t]{0.09\columnwidth}\raggedright
15\strut
\end{minipage} & \begin{minipage}[t]{0.23\columnwidth}\raggedright
table2.csv\strut
\end{minipage} & \begin{minipage}[t]{0.23\columnwidth}\raggedright
\strut
\end{minipage}\tabularnewline
\begin{minipage}[t]{0.13\columnwidth}\raggedright
Table 3\strut
\end{minipage} & \begin{minipage}[t]{0.18\columnwidth}\raggedright
02\_analysis/table2and3.do\strut
\end{minipage} & \begin{minipage}[t]{0.09\columnwidth}\raggedright
145\strut
\end{minipage} & \begin{minipage}[t]{0.23\columnwidth}\raggedright
table3.csv\strut
\end{minipage} & \begin{minipage}[t]{0.23\columnwidth}\raggedright
\strut
\end{minipage}\tabularnewline
\begin{minipage}[t]{0.13\columnwidth}\raggedright
Figure 1\strut
\end{minipage} & \begin{minipage}[t]{0.18\columnwidth}\raggedright
n.a. (no data)\strut
\end{minipage} & \begin{minipage}[t]{0.09\columnwidth}\raggedright
\strut
\end{minipage} & \begin{minipage}[t]{0.23\columnwidth}\raggedright
\strut
\end{minipage} & \begin{minipage}[t]{0.23\columnwidth}\raggedright
Source: Herodus (2011)\strut
\end{minipage}\tabularnewline
\begin{minipage}[t]{0.13\columnwidth}\raggedright
Figure 2\strut
\end{minipage} & \begin{minipage}[t]{0.18\columnwidth}\raggedright
02\_analysis/fig2.do\strut
\end{minipage} & \begin{minipage}[t]{0.09\columnwidth}\raggedright
\strut
\end{minipage} & \begin{minipage}[t]{0.23\columnwidth}\raggedright
figure2.png\strut
\end{minipage} & \begin{minipage}[t]{0.23\columnwidth}\raggedright
\strut
\end{minipage}\tabularnewline
\begin{minipage}[t]{0.13\columnwidth}\raggedright
Figure 3\strut
\end{minipage} & \begin{minipage}[t]{0.18\columnwidth}\raggedright
02\_analysis/fig3.do\strut
\end{minipage} & \begin{minipage}[t]{0.09\columnwidth}\raggedright
\strut
\end{minipage} & \begin{minipage}[t]{0.23\columnwidth}\raggedright
figure-robustness.png\strut
\end{minipage} & \begin{minipage}[t]{0.23\columnwidth}\raggedright
Requires confidential data\strut
\end{minipage}\tabularnewline
\bottomrule
\end{longtable}

\hypertarget{references}{%
\subsection{References}\label{references}}

\begin{quote}
INSTRUCTIONS: As in any scientific manuscript, you should have proper
references. For instance, in this sample README, we cited ``Ruggles et
al, 2019'' and ``DESE, 2019'' in a Data Availability Statement. The
reference should thus be listed here, in the style of your journal:
\end{quote}

Steven Ruggles, Steven M. Manson, Tracy A. Kugler, David A. Haynes II,
David C. Van Riper, and Maryia Bakhtsiyarava. 2018. ``IPUMS Terra:
Integrated Data on Population and Environment: Version 2
{[}dataset{]}.'' Minneapolis, MN: \emph{Minnesota Population Center,
IPUMS}. https://doi.org/10.18128/D090.V2

Department of Elementary and Secondary Education (DESE), 2019. ``Student
outcomes database {[}dataset{]}'' \emph{Massachusetts Department of
Elementary and Secondary Education (DESE)}. Accessed January 15, 2019.

U.S. Bureau of Economic Analysis (BEA). 2016. ``Table 30:''Economic
Profile by County, 1969-2016.'' (accessed Sept 1, 2017).

Inglehart, R., C. Haerpfer, A. Moreno, C. Welzel, K. Kizilova, J.
Diez-Medrano, M. Lagos, P. Norris, E. Ponarin \& B. Puranen et
al.~(eds.). 2014. World Values Survey: Round Six - Country-Pooled
Datafile Version:
http://www.worldvaluessurvey.org/WVSDocumentationWV6.jsp. Madrid: JD
Systems Institute.

\begin{center}\rule{0.5\linewidth}{0.5pt}\end{center}

\hypertarget{acknowledgements}{%
\subsection{Acknowledgements}\label{acknowledgements}}

Some content on this page was copied from
\href{https://www.hindawi.com/research.data/\#statement.templates}{Hindawi}.
Other content was adapted from
\href{https://doi.org/10.1093/restud/rdw057}{Fort (2016)}, Supplementary
data, with the author's permission.

\end{document}










\end{document}




